\subsection*{Construction Obligations}

\begin{remark*}[Exposition]
A number-system construction is a structured promise. The construction must
say what the objects are, how arithmetic is performed, which objects are
identified, and which algebraic or order laws the resulting system satisfies.
\end{remark*}

\begin{definitionbox}{Definition}
\begin{definition}[Number-System Construction Target]
\label{def:number-system-construction-target}
A \emph{number-system construction target} is the specified algebraic and
order structure that a proposed system is meant to satisfy. The target fixes
the required carrier, operations, relations, distinguished elements, and laws.
\end{definition}

\begin{remark*}[Standard quantified statement]
A construction target specifies the structure and laws the new arithmetic
system is intended to satisfy.
\end{remark*}

\begin{remark*}[Interpretation]
The target determines the proof obligations; without it, the phrase number
system has no fixed mathematical content.
\end{remark*}
\end{definitionbox}

\begin{dependencies}
\NoLocalDependencies
\end{dependencies}

\begin{definitionbox}{Definition}
\begin{definition}[Construction Data]
\label{def:number-system-construction-data}
The \emph{construction data} for a number system consists of the raw carrier
objects, the intended equality or equivalence relation, the distinguished
elements, and the operation formulas from which the system is built.
\end{definition}

\begin{remark*}[Standard quantified statement]
Construction data specifies the raw objects, identifications, distinguished
elements, and operation formulas used in the construction.
\end{remark*}

\begin{remark*}[Interpretation]
The data is the material being tested against the target; it is not yet a
verified number system.
\end{remark*}
\end{definitionbox}

\begin{dependencies}
\begin{itemize}
  \item \hyperref[def:number-system-construction-target]{Number-System Construction Target}
\end{itemize}
\end{dependencies}

\begin{definitionbox}{Definition}
\begin{definition}[Construction Verification]
\label{def:number-system-construction-verification}
\emph{Construction verification} is the proof that the construction data
actually satisfies the construction target. It includes, as applicable,
closure of operations, equivalence-relation proofs, well-definedness of
operations on classes, algebraic laws, order laws, and compatibility of
embeddings with earlier systems.
\end{definition}

\begin{remark*}[Standard quantified statement]
Construction verification proves that the construction data satisfies the
chosen construction target.
\end{remark*}

\begin{remark*}[Interpretation]
Verification is the step that turns proposed arithmetic into a justified
structure.
\end{remark*}
\end{definitionbox}

\begin{dependencies}
\begin{itemize}
  \item \hyperref[def:number-system-construction-target]{Number-System Construction Target}
  \item \hyperref[def:number-system-construction-data]{Construction Data}
\end{itemize}
\end{dependencies}
